\chapter{Analyse détaillée des 9 blocs de compétences RNCP38152}

Ce chapitre présente une analyse approfondie de la correspondance entre mon parcours professionnel et les 9 blocs de compétences du référentiel RNCP38152 du Master MEEF Pratiques et Ingénierie de la Formation (PIF).

\section{BC01 - Mettre en œuvre les usages avancés et spécialisés des outils numériques}

\subsection{Définition officielle du bloc}
\begin{quote}
\textit{« Identifier les usages numériques et les impacts de leur évolution sur le ou les domaines concernés par la mention. Se servir de façon autonome des outils numériques avancés pour un ou plusieurs métiers ou secteurs de recherche du domaine. »} \parencite{rncp38152}
\end{quote}

\subsection{Correspondance avec mon parcours}

\subsubsection{Expertise technique et ingénierie numérique}
Mon parcours de 20 ans dans l'industrie logicielle m'a permis de développer une expertise avancée des outils numériques, fondamentale pour l'ingénierie de formation moderne.
\begin{itemize}[leftmargin=*]
    \item \textbf{Maîtrise technologique} : Langages (Python, SQL, Web), architectures systèmes, et outils de versionning (Git/GitLab), essentiels pour concevoir des environnements d'apprentissage numériques (EIAH).
    \item \textbf{Création d'outils pour la formation} : Développement de la plateforme « Notation à l'Enveloppe » (Partie 2, Expérience 3) et du moteur de jeu éducatif (Partie 2, Expérience 2).
\end{itemize}

\subsubsection{Impact sur l'ingénierie de formation}
Je ne suis pas seulement utilisateur, mais concepteur d'outils numériques pour la formation.
\begin{itemize}[leftmargin=*]
    \item \textbf{Forge des Communs Numériques} : Contribution active au développement et partage de ressources éducatives libres (REL).
    \item \textbf{IA et Éducation} : Analyse critique et expérimentation des impacts de l'IA générative sur les pratiques de formation (Projet de thèse).
\end{itemize}

\textbf{Niveau de maîtrise : Expert} - Capacité démontrée à créer des outils numériques sur mesure pour des besoins de formation spécifiques.

\vspace{1cm}

\section{BC02 - Mobiliser et produire des savoirs hautement spécialisés}

\subsection{Définition officielle du bloc}
\begin{quote}
\textit{« Mobiliser des savoirs hautement spécialisés... Développer une conscience critique... Résoudre des problèmes... Conduire une analyse réflexive... »} \parencite{rncp38152}
\end{quote}

\subsection{Correspondance avec mon parcours}

\subsubsection{Savoirs disciplinaires et didactiques}
\begin{itemize}[leftmargin=*]
    \item \textbf{Double compétence} : Alliance de savoirs techniques pointus (NSI, développement) et de savoirs pédagogiques (didactique de l'informatique, ingénierie de formation).
    \item \textbf{Analyse de pratiques} : Utilisation de cadres théoriques (didactique professionnelle, approche par compétences) pour analyser les situations de formation (cf. Partie 2).
\end{itemize}

\subsubsection{Production de savoirs nouveaux}
\begin{itemize}[leftmargin=*]
    \item \textbf{Innovation méthodologique} : Conception de la méthode de « Notation à l'Enveloppe », une réponse originale à la problématique de l'évaluation des compétences transversales.
    \item \textbf{Recherche} : Engagement dans une démarche de recherche-action sur l'intégration de l'IA en didactique de l'informatique.
\end{itemize}

\textbf{Niveau de maîtrise : Avancé} - Capacité à articuler savoirs théoriques et savoirs d'action pour innover.

\vspace{1cm}

\section{BC03 - Mettre en œuvre une communication spécialisée pour le transfert de connaissances}

\subsection{Définition officielle du bloc}
\begin{quote}
\textit{« Identifier, sélectionner et analyser avec esprit critique diverses ressources spécialisées. Communiquer à des fins de formation ou de transfert de connaissances... »} \parencite{rncp38152}
\end{quote}

\subsection{Correspondance avec mon parcours}

\subsubsection{Communication pédagogique et formation d'adultes}
\begin{itemize}[leftmargin=*]
    \item \textbf{Formation de formateurs} : Animation d'ateliers et accompagnement de collègues sur les outils numériques et la pédagogie de projet.
    \item \textbf{Vulgarisation technique} : Capacité à rendre accessibles des concepts informatiques complexes (SQL, Algorithmique) à des publics variés (élèves, étudiants IUT, collègues).
\end{itemize}

\subsubsection{Médiation scientifique et technique}
\begin{itemize}[leftmargin=*]
    \item \textbf{Documentation et partage} : Rédaction de documentations techniques pédagogiques sur la Forge et le site \textit{home.educ-ai.fr}.
    \item \textbf{Rayonnement} : Communication professionnelle lors de groupes de travail académiques.
\end{itemize}

\textbf{Niveau de maîtrise : Avancé} - Expertise dans la transmission de savoirs techniques et méthodologiques.

\vspace{1cm}

\section{BC04 - Contribuer à la transformation en contexte professionnel}

\subsection{Définition officielle du bloc}
\begin{quote}
\textit{« Gérer des contextes professionnels complexes... Conduire un projet... Analyser ses actions... Respecter les principes d'éthique... »} \parencite{rncp38152}
\end{quote}

\subsection{Correspondance avec mon parcours}

\subsubsection{Pilotage de projets complexes}
\begin{itemize}[leftmargin=*]
    \item \textbf{Gestion de projet industriel} : 12 ans d'expérience en tant que Team Leader, gestion de budgets et d'équipes, transférable à la gestion de projets de formation.
    \item \textbf{Projets interdisciplinaires} : Pilotage du projet « Jeu Vidéo » (Partie 2), coordonnant 5 enseignants et 40 élèves sur 8 mois.
\end{itemize}

\subsubsection{Démarche qualité et éthique}
\begin{itemize}[leftmargin=*]
    \item \textbf{Amélioration continue} : Application des méthodes agiles (Scrum, feedback) à l'ingénierie de formation.
    \item \textbf{Éthique du numérique} : Intégration systématique des enjeux RGPD, accessibilité (W3C) et inclusion dans les projets numériques développés.
\end{itemize}

\textbf{Niveau de maîtrise : Expert} - Compétences éprouvées en gestion de projet et accompagnement du changement.

\vspace{1cm}

\section{BC05 - Concevoir, piloter et animer des dispositifs ou ingénieries de formation ou de médiation}

\subsection{Définition officielle du bloc}
\begin{quote}
\textit{« Recueillir des besoins, formaliser un projet de formation. Établir et mettre en œuvre un cahier des charges. Mobiliser des savoirs spécialisés... »} \parencite{rncp38152}
\end{quote}

\subsection{Correspondance avec mon parcours}

\subsubsection{Ingénierie didactique et pédagogique}
\begin{itemize}[leftmargin=*]
    \item \textbf{Analyse des besoins} : Diagnostic des difficultés d'apprentissage du SQL (Partie 2) pour concevoir un dispositif de remédiation par rétro-ingénierie.
    \item \textbf{Conception de dispositifs} : Création de scénarios pédagogiques complets (progressions, séquences, évaluations) adaptés aux publics (Lycée, IUT).
\end{itemize}

\subsubsection{Adaptation aux publics}
\begin{itemize}[leftmargin=*]
    \item \textbf{Différenciation} : Conception de parcours adaptés aux élèves à besoins particuliers (ULIS, troubles dys) grâce au numérique.
    \item \textbf{Publics hétérogènes} : Adaptation des stratégies pédagogiques pour des groupes variés (scientifiques, tertiaires).
\end{itemize}

\textbf{Niveau de maîtrise : Avancé} - Capacité à concevoir des ingénieries complètes, de l'analyse des besoins à l'évaluation.

\vspace{1cm}

\section{BC06 - Situer ses interventions dans un écosystème professionnel, coopérer, co-construire des actions}

\subsection{Définition officielle du bloc}
\begin{quote}
\textit{« Inscrire son action dans un secteur professionnel de formation... Conduire un projet dans un cadre collaboratif... Contribuer à la performance stratégique d'une équipe. »} \parencite{rncp38152}
\end{quote}

\subsection{Correspondance avec mon parcours}

\subsubsection{Travail collaboratif et réseau}
\begin{itemize}[leftmargin=*]
    \item \textbf{Co-construction} : Travail en équipe pluridisciplinaire (Lettres, Arts, Langues, Info) pour le projet Jeu Vidéo.
    \item \textbf{Réseaux professionnels} : Participation active aux réseaux d'enseignants NSI et à la communauté de la Forge des Communs Numériques.
\end{itemize}

\subsubsection{Partenariats}
\begin{itemize}[leftmargin=*]
    \item \textbf{Lien École-Entreprise} : Mobilisation de mon réseau professionnel pour l'orientation des élèves et l'intervention d'experts.
    \item \textbf{Écosystème numérique} : Collaboration avec les acteurs du numérique éducatif (DNE, académie).
\end{itemize}

\textbf{Niveau de maîtrise : Acquis} - Forte culture de la coopération issue du monde de l'entreprise et réinvestie dans l'éducation.

\vspace{1cm}

\section{BC07 - Produire et diffuser des ressources, des outils, des démarches, et des contenus de formation innovants}

\subsection{Définition officielle du bloc}
\begin{quote}
\textit{« Identifier des besoins en termes de ressources. Adapter la production... Faire évoluer les modalités d'enseignement... Identifier le processus de production, de diffusion et de valorisation des savoirs. »} \parencite{rncp38152}
\end{quote}

\subsection{Correspondance avec mon parcours}

\subsubsection{Conception de ressources innovantes}
\begin{itemize}[leftmargin=*]
    \item \textbf{Outils numériques} : Développement d'applications web pour l'apprentissage (Quiz, IDE en ligne, outils de gestion de projet).
    \item \textbf{Supports pédagogiques} : Création de capsules vidéo, de tutoriels interactifs et de supports de cours différenciés.
\end{itemize}

\subsubsection{Diffusion et valorisation}
\begin{itemize}[leftmargin=*]
    \item \textbf{Open Source} : Publication des outils développés sous licences libres pour favoriser leur réutilisation par la communauté éducative.
    \item \textbf{Partage d'expertise} : Diffusion de pratiques innovantes (méthode d'évaluation, pédagogie de projet) auprès des pairs.
\end{itemize}

\textbf{Niveau de maîtrise : Expert} - Production significative de ressources numériques utilisées par d'autres formateurs.

\vspace{1cm}

\section{BC08 - Conduire, piloter et partager une analyse réflexive sur les pratiques}

\subsection{Définition officielle du bloc}
\begin{quote}
\textit{« Accompagner des professionnels en formation dans l'analyse de leurs pratiques... Former les formateurs à prendre en compte les critiques et à réajuster les pratiques. »} \parencite{rncp38152}
\end{quote}

\subsection{Correspondance avec mon parcours}

\subsubsection{Analyse de pratiques}
\begin{itemize}[leftmargin=*]
    \item \textbf{Auto-analyse} : Pratique réflexive systématique sur mes propres enseignements (cf. bilans d'expériences Partie 2).
    \item \textbf{Accompagnement} : Rôle de tuteur informel auprès de collègues débutants en NSI, aide à l'analyse de leurs difficultés techniques et pédagogiques.
\end{itemize}

\subsubsection{Démarche qualité}
\begin{itemize}[leftmargin=*]
    \item \textbf{Évaluation des dispositifs} : Mise en place d'indicateurs pour mesurer l'efficacité des formations et des outils (réussite aux examens, satisfaction, engagement).
    \item \textbf{Régulation} : Ajustement continu des scénarios pédagogiques basé sur l'analyse des résultats et les retours des apprenants.
\end{itemize}

\textbf{Niveau de maîtrise : Avancé} - Posture réflexive ancrée et capacité à accompagner l'analyse de pratiques.

\vspace{1cm}

\section{BC09 - S'engager dans une démarche de développement professionnel}

\subsection{Définition officielle du bloc}
\begin{quote}
\textit{« Se tenir informé des évolutions de la recherche scientifique... Conduire une analyse réflexive... Développer des ingénieries pédagogiques et didactiques dans une perspective d'enseignement, de formation, de recherche. »} \parencite{rncp38152}
\end{quote}

\subsection{Correspondance avec mon parcours}

\subsubsection{Développement professionnel continu}
\begin{itemize}[leftmargin=*]
    \item \textbf{Formation continue} : Suivi de nombreuses formations académiques et auto-formation technique permanente.
    \item \textbf{Recherche} : Engagement dans un projet de thèse, lecture de la littérature scientifique en didactique et sciences de l'éducation.
\end{itemize}

\subsubsection{Perspective de recherche}
\begin{itemize}[leftmargin=*]
    \item \textbf{Ingénierie didactique} : Conception de situations d'apprentissage fondées sur la recherche (problématisation, institutionnalisation).
    \item \textbf{Veille} : Suivi actif des innovations en technologies éducatives (IA, Learning Analytics).
\end{itemize}

\textbf{Niveau de maîtrise : Acquis} - Engagement fort dans une dynamique de chercheur-praticien.
